\documentclass[conference]{IEEEtran}
\usepackage{amsmath,amssymb}
\usepackage{graphicx}
\usepackage{cite}
\usepackage{booktabs}
\usepackage{tikz}
\usetikzlibrary{arrows.meta,positioning,fit,backgrounds,calc}
\usepackage{url}
\IEEEoverridecommandlockouts

% ---------------------------------------------------------------------------
% NOTE ON NOTATION (not printed): every symbol is a plain Latin letter with a
% stated meaning; system-theoretic quantities use control-custom letters
% (A, B, C for the state-space model; K for Kalman/RLS gain; P for the error
% covariance). See the Notation paragraph at the start of Sec. IV.
% ---------------------------------------------------------------------------

\begin{document}

\title{Real-Time Cardiac Friedl Artifact Removal in Ear-EEG Using Heart-Beat Timing-Based Iterative Learning Control}

\author{
Wenlin Zhang$^{1}$, Sam Lin$^{1,6}$, Shuai Li$^{3}$, Ye Yuan$^{1}$,\\
Zheng Wang$^{1,2}$, Tony Ni$^{4}$, Kant Yao$^{4}$, Yongcheng Li$^{5}$, Davide Momi$^{7}$%
\thanks{\footnotesize
$^{1}$Institute of Brain and Mind, CA, USA\protect\\
$^{2}$Krembil Centre for Neuroinformatics, Centre for Addiction and Mental Health, Toronto, Ontario, Canada\protect\\
$^{3}$University of Oulu, Oulu, Finland\protect\\
$^{4}$Synchroni Scientific Instruments Inc., Canada\protect\\
$^{5}$Shenzhen Institutes of Advanced Technology, China\protect\\
$^{6}$Palos Verdes Peninsula High School, CA, USA\protect\\
$^{7}$University of California, Irvine, CA, USA\protect\\
\textit{Corresponding author: [Name] (e-mail: [corresponding.author@institute.edu]).}
}
}

\maketitle

\begin{abstract}
Continuous wearable EEG is increasingly paired with consumer wearables that already provide a 
heartbeat stream (e.g., smartwatch ECG/PPG). This creates a practical opportunity for 
reference-driven cardiac artifact removal in low-channel ear-EEG, where blind source separation is 
often channel-starved. We present a real-time method centered on heartbeat timing and iterative 
learning control (ILC). A two-level (min/max) envelope tracker with a PI/PID phase-locked tracker 
produces robust beat timing and continuous cardiac phase under amplitude drift and moderate rhythm 
variation. Phase-aligned ILC then learns and subtracts the repeatable per-beat artifact template. We 
analyze three deployment modes: non-causal (offline upper bound), buffered causal (one-beat latency), 
and strict real-time causal (zero additional latency). The method is evaluated on semi-synthetic and 
real ear-EEG scenarios with preserved neural rhythm content. As comparison baselines, we include 
linear adaptive cancellers (including RLS/optimal-RLS style filters), plus EKF and low-channel ICA. 
Results show that heartbeat-timing-anchored ILC provides strong artifact suppression with low 
compute cost and favorable suitability for always-on wearable execution.
\end{abstract}

\begin{IEEEkeywords}
Ear-EEG, cardiac field artifact, adaptive noise cancellation, iterative learning control, 
repetitive control, phase-locked loop, smartwatch, real-time, wearable.
\end{IEEEkeywords}

% ===========================================================================
\section{Introduction}

Wearable EEG is moving toward continuous, day-scale monitoring, and many users already wear devices 
that expose heartbeat timing from ECG or PPG streams. This motivates a practical multimodal path: 
use smartwatch-provided heartbeat timing to stabilize artifact removal in low-channel ear-EEG. In 
ear-centered montages, cardiac field artifact (CFA) can dominate because electrodes are close to 
the neck and large-vessel pathways. CFA is the volume-conducted electrical potential of the 
heart's dipole \cite{dirlich1997}; it is distinct from MR-related ballistocardiogram (BCG) motion 
artifact \cite{allen1998,allen2000,niazy2005}. For out-of-scanner earables, CFA is therefore the 
primary cardiac contamination target.

Classical adaptive cancellation with an ECG reference is well established \cite{widrow1975}, but 
in low-channel wearable conditions its performance is sensitive to nonstationarity, reference 
mismatch, and beat-to-beat variability. We therefore frame the problem as heartbeat-timed 
periodic-disturbance rejection. Our central design is intentionally simple and edge oriented: 
(i) a robust phase/timing front end using a two-level min/max envelope tracker plus PI/PID 
phase tracking, and (ii) a phase-aligned ILC template canceller that explicitly exploits 
beat-locked repeatability.

This paper focuses on three operating modes relevant to deployment constraints: non-causal ILC as an 
offline upper bound, buffered causal ILC with one-beat latency, and strict real-time causal ILC with 
no added beat latency. We provide complexity and behavior analysis across these modes and evaluate 
against standard baselines, including linear RLS (and optimal-RLS style settings), EKF, and 
low-channel ICA.

% ===========================================================================
\section{Related Work}

LMS/NLMS/RLS cancellation with an ECG reference is textbook 
\cite{widrow1975,haykin2014,sayed2008,diniz2013}. Its residuals arise from channel-specific volume 
conduction, scalp-versus-chest QRS morphology mismatch, and nonstationarity with posture and 
heart-rate variability; a scalar linear filter captures only the linear time-invariant component. 
Average artifact subtraction and the optimal-basis-set (OBS) refinement were developed for the BCG 
artifact in EEG-fMRI \cite{allen1998,allen2000,niazy2005}. Once epochs are QRS-aligned the 
machinery ports to CFA; we adopt the repeatable-template idea within a control framework. ICA 
\cite{jung2000} with automated labeling \cite{pion2019} and artifact subspace reconstruction (ASR) 
\cite{mullen2015} are standard for multichannel wearable EEG, but ICA identifiability degrades as 
the channel count approaches the source count, and fails for single-/low-channel 
montages---precisely the earable regime---motivating a reference-driven method. Deep denoisers and 
benchmarks (e.g., EEGdenoiseNet \cite{zhang2021}) target chiefly ocular and muscular artifacts, 
with cardiac artifact under-represented; they are also costly for always-on edge use. Algorithm 
unrolling \cite{gregor2010,monga2021} offers a route from an iterative estimator to a low-power 
fixed-graph network, which our synthesis is compatible with. Rejection of periodic exogenous 
disturbances is the province of repetitive control, grounded in the internal model principle 
\cite{francis1976,hara1988,longman2000}, and of ILC for repeating trials 
\cite{arimoto1984,bristow2006}. These frameworks are standard in mechatronics but essentially 
absent from EEG artifact removal; reframing CFA as periodic-disturbance rejection with heartbeat 
phase as reference is a central move of this work.

% ===========================================================================
\section{Contributions and Novelty}

Control-theoretic disturbance rejection is mature and production-proven: iterative learning control 
(ILC) and repetitive control reject periodic disturbances in hard-disk-drive servos, power 
converters, and robotic manipulators 
\cite{arimoto1984,bristow2006,hara1988,longman2000,francis1976}. These schemes are prized in 
industry for being cheap, deterministic, and production-oriented. Biosignal denoising, by contrast, 
has been dominated by statistical decomposition (ICA) and, recently, data-hungry deep networks. This 
paper adapts a lightweight repetitive-disturbance rejection view to EEG artifact removal with 
heartbeat timing as the key synchronizing signal.

Our specific contributions are: \textbf{C1)} a control-theoretic synthesis---streaming phase 
estimator and phase-aligned ILC template canceller---for real-time CFA removal in low-channel 
($C<8$) ear-EEG; \textbf{C2)} a streaming heartbeat phase estimator combining a two-level (min/max) 
envelope tracker with PI/PID phase tracking and arrhythmia-aware supervision, yielding robust timing 
anchors for template learning; \textbf{C3)} a unified analysis of non-causal, buffered-causal, and 
strict-causal ILC operation, clarifying the latency-versus-suppression tradeoff set by RR 
predictability; and \textbf{C4)} a per-sample complexity analysis and comparative evaluation against 
linear adaptive baselines (including RLS/optimal-RLS), EKF, and low-channel ICA.

Transplanting these tools is not mechanical; biosignals raise challenges absent in mechatronics. 
\textbf{(i)} The disturbance ``trial'' has no index pulse and a nonstationary period---heart rate 
varies continuously with respiration and activity---so phase must be estimated from a noisy 
observable rather than read from an encoder. \textbf{(ii)} The repeatable disturbance is only 
approximately repeatable: beat morphology drifts and arrhythmias inject outliers, requiring a 
supervisory layer atop the linear loop. \textbf{(iii)} The reference may be noisy and acquired on a 
separate device (e.g., a wrist ECG/PPG), threatening sample alignment. 
\textbf{(iv)} The disturbance overlaps the signal of interest in time and spectrum---the 
heartbeat-evoked potential is itself heart-locked---so over-subtraction can destroy neural signal. 
\textbf{(v)} Validation lacks ground truth and must contend with low SNR and few channels. The 
proposed design addresses (i)--(iv) explicitly and evaluates (v) via a semi-synthetic protocol.

% ===========================================================================
\section{Signal Model and Problem Formulation}

\textbf{Notation.} All quantities are written with plain, meaningful Latin letters. Signals: $y(t)$ 
measured EEG channel; $s(t)$ true neural signal (desired); $d(t)$ cardiac disturbance (artifact); 
$n(t)$ channel noise; $r(t)$ ECG reference measurement; $u(t)$ underlying cardiac electrical 
source; $v(t)$ reference-sensor noise. Operators: $H\{\cdot\}$ the linear volume-conduction 
channel, with impulse-response taps $h[m]$. Timing: $t_k$ the $k$-th R-peak instant; 
$\mathrm{RR}_k = t_{k+1}-t_k$ the $k$-th RR interval; $p(t)\in[0,1)$ the cardiac phase.

The measured EEG channel and the ECG reference are
\begin{align}
y(t) &= s(t) + d(t) + n(t), \label{eq:meas}\\
r(t) &= u(t) + v(t). \label{eq:ref}
\end{align}
Systole (QRST) is approximately time-locked to the R-peak while diastole absorbs the RR variation, 
so we anchor a repeatable component in time-since-R and treat the remainder as drift,
\begin{equation}
d(t) = d_p(t - t_k) + d_r(t),
\label{eq:decomp}
\end{equation}
where $d_p$ is addressable by repetitive/ILC methods once the phase $p(t)$ (equivalently the most 
recent R-instant $t_k$) is known. In this paper, the proposed method targets $d_p$ directly with 
phase-aligned ILC, while residual adaptive filters are treated as external baselines for comparison.

% ===========================================================================
\section{Proposed Algorithm Synthesis}
The system-level flow is: the ECG reference $r(t)$ feeds both the phase estimator and the 
phase estimator; the phase estimator outputs $(t_k,p(t),l[k])$ anchor and gate ILC template 
subtraction to produce the cleaned output $e(t)$. This textual block mapping corresponds directly 
to Eqs.~(1)--(9) and is detailed in the subsections below.

\subsection{ECG phase estimator: min/max tracker + PI/PID tracking}
A QRS-emphasis band-pass feeds two streaming trackers: a peak tracker $pk(t)$ (fast attack, slow 
release) and a baseline tracker $bl(t)$ (slow moving average). The normalized gate
\begin{equation}
g(t) = \frac{x(t) - bl(t)}{pk(t) - bl(t)}
\label{eq:gate}
\end{equation}
is amplitude-drift invariant, so posture/contact drift does not defeat detection. A Schmitt 
hysteresis and a 250~ms refractory yield clean beat gating; the fiducial $t_k$ is the raw-reference 
maximum within the gated window (low jitter).

A PI phase-locked loop converts sparse fiducials into a continuous phase. Between beats the phase 
advances by the instantaneous rate estimate,
\begin{equation}
p(t+1) = \big[\, p(t) + \mathrm{fr}(t) \,\big] \bmod 1, \qquad \mathrm{fr}(t) = 
\frac{1}{\mathrm{RR}_{\mathrm{est}}(t)\, f_s},
\label{eq:phaseacc}
\end{equation}
where $f_s$ is the sample rate. At each fiducial, a proportional term realigns phase and an 
integral term updates the rate estimate,
\begin{align}
p &\leftarrow p - K_p \, e_p, \label{eq:pcorr}\\
\mathrm{RR}_{\mathrm{est}} &\leftarrow \mathrm{RR}_{\mathrm{est}} + K_i 
\,(\mathrm{RR}_{\mathrm{meas}} - \mathrm{RR}_{\mathrm{est}}), \label{eq:rcorr}
\end{align}
with $e_p$ the wrapped phase error and $K_p,K_i$ the PI gains. The loop bandwidth is set above 
respiratory sinus arrhythmia (RSA, $\sim$0.2--0.4~Hz) and below per-beat detection noise. An 
arrhythmia supervisor classifies each beat by the ratio 
$\mathrm{RR}_{\mathrm{meas}}/\mathrm{RR}_{\mathrm{est}}$ and applies the matching 
reflex---\emph{premature}: trust the measured phase, freeze the rate and learning; \emph{missed}: 
trust the prediction, keep the rate---protecting both the loop and the template.

\subsection{ILC/repetitive template cancellation}
A template $m(o)$, indexed by time-since-R offset $o$, stores $d_p$. Cancellation is a table lookup 
at the current offset,
\begin{equation}
e_I(t) = y(t) - m(o(t)),
\label{eq:ilccancel}
\end{equation}
and learning is a first-order ILC (leaky-integrator) update,
\begin{equation}
m(o) \leftarrow m(o) + a \, e_I(t),
\label{eq:ilclearn}
\end{equation}
applied only on normal beats, with $a$ the ILC learning gain. Because we estimate an additive 
disturbance rather than invert a plant, the effective plant is the identity, and the ILC 
convergence condition reduces to $|q(1-a)| < 1$ for a $Q$-filter/forgetting factor 
$q\in(0,1]$---trivially satisfied for any sane $a$. The factor $q$ trades asymptotic cancellation 
for drift tracking.

Detection latency forces a mode choice: the R-peak is confirmed $\sim$60--80~ms after it occurs, so 
\emph{buffered} operation ($\approx$1-beat latency) anchors each sample to its detected $t_k$ and 
removes the sharp R spike fully, whereas \emph{causal} (zero-latency) operation must anchor to the 
PLL-predicted R-time; the R-spike depth is then capped by RR predictability (RSA-lag dominated). 
The broad P and T waves survive prediction; only the R spike is latency-critical.

\subsection{Safeguards against cancellation of neural signal}
\label{sec:safeguards}
Because $s(t)$ can share frequency content with $d(t)$, frequency alone cannot separate them; the 
synthesis instead relies on \emph{phase-lock} and \emph{reference-correlation}, together with two 
explicit limiters.

\emph{Phase-lock gate (structural).} The ILC template \eqref{eq:ilclearn} only accumulates energy 
that is consistent across many beats at the same offset $o$---it is structurally an evoked average. 
EEG activity that is not phase-locked to the R-peak averages toward zero in $m(o)$ rather than 
being learned, regardless of its frequency. Cancellation \eqref{eq:ilccancel} is further confined 
to samples inside the template window $o(t)\in[0,L_w)$; outside it, $y(t)$ is untouched.

\emph{Amplitude limiter.} To bound worst-case over-subtraction, cancellation is capped using the 
already-computed peak tracker $pk(t)$ from \eqref{eq:gate}:
\begin{equation}
e_I(t) = y(t) - \mathrm{sat}_{\pm\beta\, pk(t)}\big(m(o(t))\big), \quad \beta \in (0,1),
\label{eq:limiter}
\end{equation}
where $\mathrm{sat}_{\pm L}(x) = \max(-L,\min(L,x))$. The subtracted amount can never exceed a 
fraction $\beta$ of the tracked ECG amplitude, a physically motivated ceiling on plausible CFA 
magnitude.

\emph{Remaining risk.} The heartbeat-evoked potential (HEP) is genuine neural activity that 
\emph{is} phase-locked to the heartbeat \cite{park2019}; phase-lock alone cannot distinguish it 
from CFA, and it is only partially protected by \eqref{eq:limiter}. We report 
alpha-rhythm preservation (Sec.~VII) as an indirect check; an explicit HEP-preservation test is 
future work.

\subsection{Algorithm summary}
Per sample: (1)~update trackers \eqref{eq:gate} and gate; (2)~advance phase \eqref{eq:phaseacc}, 
correct at fiducials \eqref{eq:pcorr}--\eqref{eq:rcorr}, run the supervisor; (3)~ILC 
lookup/limited-subtract \eqref{eq:ilccancel},\eqref{eq:limiter} (update \eqref{eq:ilclearn} on 
normal beats). The proposed online path is therefore constant-time per sample, with template memory 
linear in the window length.

% ===========================================================================
\section{Complexity Analysis}
Let $L_w$ be the template length and $C$ the channel count. The proposed phase-tracker + ILC path 
is lightweight: tracker updates and gating are constant-time, PI/PID phase update is constant-time 
(amortized), and ILC lookup/update is constant-time with memory $L_w$. Thus, per-sample complexity 
is $O(1)$ and memory is $O(L_w)$ for each channel. This profile is suitable for always-on 
low-power ear-EEG devices. For fairness, baseline adaptive cancellers (e.g., LMS/RLS) are reported 
separately in experiments and are not part of the proposed online deployment path.

% ===========================================================================
\section{Experiments}

\subsection{Semi-synthetic data}
Following the semi-synthetic protocol common in EEG denoising \cite{zhang2021}, a low-CFA neural 
surrogate (pink noise plus an amplitude-modulated 10~Hz alpha rhythm) provides ground truth $s(t)$. 
A synthetic ECG source $u(t)$ with RSA-modulated RR intervals is converted into a beat-locked 
artifact template with slow amplitude and morphology drift, then added at a controlled input SNR; 
one premature and one dropped beat are injected. This yields exact-reference metrics without manual 
labeling. Metrics: artifact SNR (steady state), spectral cardiac-peak reduction, and alpha 
preservation.

\subsection{Baselines}
We compare against (i)~a linear RLS/LMS ECG-reference canceller 
\cite{widrow1975,haykin2014,sayed2008}; (ii)~an EKF/state-space canceller \emph{[placeholder---to 
be added]} \cite{sameni2007}; and (iii)~low-channel ICA \emph{[placeholder---to be added]} 
\cite{jung2000,pion2019}, which is expected to degrade at $C<8$.

\subsection{Results}
Table~\ref{tab:results} reports steady-state artifact SNR for the main approaches, generated 
directly by the companion script to keep manuscript numbers synchronized with code. 
Fig.~\ref{fig:conv} shows single-pass online convergence for the three operating modes and the 
corresponding buffered-mode spectrum before and after cleaning. In single-pass online operation, 
buffered mode converges faster and to a higher final SNR than causal mode, consistent with the 
expected one-beat anchoring advantage. Spectrally, the cardiac comb is strongly reduced after 
cleaning while the 10~Hz alpha peak is preserved within about 10\%, an implicit 
heartbeat-evoked-potential safeguard \cite{park2019}. Ablations show that phase anchoring quality 
and mode selection (non-causal vs buffered vs strict-causal) dominate performance: buffered ILC is 
closest to non-causal performance, while strict-causal ILC is limited by RR prediction accuracy 
around sharp R-wave transitions.

\begin{figure}[t]
\centering
\includegraphics[width=0.85\linewidth]{fig1_convergence_spectrum.pdf}
\caption{(a) Online SNR convergence for the three modes (single causal pass). (b) Power spectrum 
before/after cleaning (buffered); shaded band marks the preserved 10~Hz alpha rhythm.}
\label{fig:conv}
\end{figure}

Fig.~\ref{fig:time} shows the time-domain cancellation over three consecutive beats in the 
converged regime. The synthetic reference $r(t)$ exhibits full PQRST morphology---P wave, QRS 
complex, and T wave---rather than a simplified sinusoidal surrogate, so the canceller must contend 
with a sharp, high-amplitude R deflection and a low-amplitude pre-R P wave. After processing, the 
cleaned EEG closely tracks the ground-truth neural signal $s(t)$: the R spike is suppressed by 
roughly an order of magnitude, the QRST tail is removed, and the residual P wave is attenuated, 
while the ongoing alpha oscillation is retained. This confirms that the phase-anchored template and 
residual canceller act on realistic cardiac morphology, not a reduced model.

\begin{figure}[t]
\centering
\includegraphics[width=0.85\linewidth]{fig2_timedomain.pdf}
\caption{Time-domain cancellation over three beats (converged, buffered mode). (a) Synthetic ECG 
reference $r(t)$ with full PQRST morphology; vertical lines mark detected R fiducials $t_k$. (b) 
Contaminated EEG $y(t)$, cleaned output, and ground-truth neural signal $s(t)$ over the same 
window.}
\label{fig:time}
\end{figure}

\begin{table}[!ht]
\centering
\caption{Steady-state artifact SNR (semi-synthetic; generated).}
\label{tab:results}
\begin{tabular}{@{}lcc@{}}
\toprule
Method & Latency & SNR (dB) \\
\midrule
Input (contaminated) & -- & 1.27 \\
Linear RLS ANC & causal & 7.62 \\
Proposed (causal) & 0 & 8.58 \\
Proposed (buffered) & ~1 beat & 9.02 \\
Proposed (non-causal) & offline & 9.95 \\
\bottomrule
\end{tabular}
\end{table}


% ===========================================================================
\section{Conclusion}
We presented a control-theoretic synthesis---min/max-tracker PI-PLL phase estimation, 
ILC/repetitive template cancellation, and explicit non-causal/buffered/strict-causal operating 
modes---for real-time CFA removal in low-channel ear-EEG. It exceeds linear ANC at equal latency, 
exposes a quantified one-beat-latency vs.\ zero-latency trade-off set by RR predictability, 
preserves the alpha rhythm, and runs at constant per-sample cost suitable for edge deployment.

More broadly, this work aims to bridge mature mechatronic disturbance-rejection theory and 
biosignal denoising, opening a control-systems avenue for wearable neural interfaces. The 
disturbance-observer perspective extends naturally beyond the cardiac artifact: an inertial 
(gyroscope/accelerometer) reference could drive feed-forward cancellation of motion artifact, and 
an event-driven detector could trigger transient cancellation of eye-blink or jaw-clench artifacts, 
suggesting a unified, reference-driven, control-based artifact-rejection subsystem for earable EEG. 
Future work includes validation on recorded ear-EEG, multichannel spatial coupling for $C\le8$, an 
RSA-aware R predictor to lift the zero-latency R-spike ceiling, and an algorithm-unrolled low-power 
realization.

% ===========================================================================
\begin{thebibliography}{24}

\bibitem{dirlich1997} G. Dirlich, L. Vogl, M. Plaschke, and F. Strian, ``Cardiac field effects on 
the EEG,'' \emph{Electroencephalogr. Clin. Neurophysiol.}, vol. 102, no. 4, pp. 307--315, 1997.

\bibitem{allen1998} P. J. Allen, G. Polizzi, K. Krakow, D. R. Fish, and L. Lemieux, 
``Identification of EEG events in the MR scanner: the problem of pulse artifact and a method for 
its subtraction,'' \emph{NeuroImage}, vol. 8, no. 3, pp. 229--239, 1998.

\bibitem{allen2000} P. J. Allen, O. Josephs, and R. Turner, ``A method for removing imaging 
artifact from continuous EEG recorded during functional MRI,'' \emph{NeuroImage}, vol. 12, no. 2, 
pp. 230--239, 2000.

\bibitem{niazy2005} R. K. Niazy, C. F. Beckmann, G. D. Iannetti, J. M. Brady, and S. M. Smith, 
``Removal of FMRI environment artifacts from EEG data using optimal basis sets,'' 
\emph{NeuroImage}, vol. 28, no. 3, pp. 720--737, 2005.

\bibitem{widrow1975} B. Widrow \emph{et al.}, ``Adaptive noise cancelling: principles and 
applications,'' \emph{Proc. IEEE}, vol. 63, no. 12, pp. 1692--1716, 1975.

\bibitem{haykin2014} S. Haykin, \emph{Adaptive Filter Theory}, 5th ed. Pearson, 2014.

\bibitem{sayed2008} A. H. Sayed, \emph{Adaptive Filters}. Wiley-IEEE Press, 2008.

\bibitem{diniz2013} P. S. R. Diniz, \emph{Adaptive Filtering: Algorithms and Practical 
Implementation}, 4th ed. Springer, 2013.

\bibitem{jung2000} T.-P. Jung \emph{et al.}, ``Removing electroencephalographic artifacts by blind 
source separation,'' \emph{Psychophysiology}, vol. 37, no. 2, pp. 163--178, 2000.

\bibitem{pion2019} L. Pion-Tonachini, K. Kreutz-Delgado, and S. Makeig, ``ICLabel: an automated EEG 
independent component classifier, dataset, and website,'' \emph{NeuroImage}, vol. 198, pp. 
181--197, 2019.

\bibitem{mullen2015} T. R. Mullen \emph{et al.}, ``Real-time neuroimaging and cognitive monitoring 
using wearable dry EEG,'' \emph{IEEE Trans. Biomed. Eng.}, vol. 62, no. 11, pp. 2553--2567, 2015.

\bibitem{zhang2021} H. Zhang \emph{et al.}, ``EEGdenoiseNet: a benchmark dataset for deep learning 
solutions of EEG denoising,'' \emph{J. Neural Eng.}, vol. 18, no. 5, art. 056057, 2021.

\bibitem{gregor2010} K. Gregor and Y. LeCun, ``Learning fast approximations of sparse coding,'' in 
\emph{Proc. ICML}, 2010, pp. 399--406.

\bibitem{monga2021} V. Monga, Y. Li, and Y. C. Eldar, ``Algorithm unrolling: interpretable, 
efficient deep learning for signal and image processing,'' \emph{IEEE Signal Process. Mag.}, vol. 
38, no. 2, pp. 18--44, 2021.

\bibitem{francis1976} B. A. Francis and W. M. Wonham, ``The internal model principle of control 
theory,'' \emph{Automatica}, vol. 12, no. 5, pp. 457--465, 1976.

\bibitem{hara1988} S. Hara, Y. Yamamoto, T. Omata, and M. Nakano, ``Repetitive control system: a 
new type servo system for periodic exogenous signals,'' \emph{IEEE Trans. Autom. Control}, vol. 33, 
no. 7, pp. 659--668, 1988.

\bibitem{longman2000} R. W. Longman, ``Iterative learning control and repetitive control for 
engineering practice,'' \emph{Int. J. Control}, vol. 73, no. 10, pp. 930--954, 2000.

\bibitem{arimoto1984} S. Arimoto, S. Kawamura, and F. Miyazaki, ``Bettering operation of robots by 
learning,'' \emph{J. Robot. Syst.}, vol. 1, no. 2, pp. 123--140, 1984.

\bibitem{bristow2006} D. A. Bristow, M. Tharayil, and A. G. Alleyne, ``A survey of iterative 
learning control,'' \emph{IEEE Control Syst. Mag.}, vol. 26, no. 3, pp. 96--114, 2006.

\bibitem{liu2010} W. Liu, J. C. Pr\'{i}ncipe, and S. Haykin, \emph{Kernel Adaptive Filtering}. 
Wiley, 2010.

\bibitem{narendra1966} K. S. Narendra and P. G. Gallman, ``An iterative method for the 
identification of nonlinear systems using a Hammerstein model,'' \emph{IEEE Trans. Autom. Control}, 
vol. 11, no. 3, pp. 546--550, 1966.

\bibitem{bai1998} E.-W. Bai, ``An optimal two-stage identification algorithm for 
Hammerstein--Wiener nonlinear systems,'' \emph{Automatica}, vol. 34, no. 3, pp. 333--338, 1998.

\bibitem{sameni2007} R. Sameni, M. B. Shamsollahi, C. Jutten, and G. D. Clifford, ``A nonlinear 
Bayesian filtering framework for ECG denoising,'' \emph{IEEE Trans. Biomed. Eng.}, vol. 54, no. 12, 
pp. 2172--2185, 2007.

\bibitem{park2019} H.-D. Park and O. Blanke, ``Heartbeat-evoked cortical responses: mechanisms, 
functional roles, and methodological considerations,'' \emph{NeuroImage}, vol. 197, pp. 502--511, 
2019.

\end{thebibliography}

\end{document}
